\chapter{Systematic Literature Review Protocol}
\chaptermark{SLR Protocol}{}

This chapter presents the protocol for the Systematic Literature Review
(\acronym{SLR}) which serves as the foundation for the research component of
this thesis. The protocol has been developed in accordance with the
\acronym{PRISMA} guidelines \cite{page2021prisma} and the \acronym{PICOC}
framework \cite{kitchenham2007guidelines}.

\section{Review Identification}

The objective of this review is to synthesize current methodologies that
utilize \termdef{Meta-Optimization} to dynamically schedule loss functions
representing decomposed sub-tasks of varying difficulty (e.g., via function
approximation, smoothing, or auxiliary objectives) in deep learning scenarios.

\section{The PICOC Framework}

To precisely define the scope of the search, the \acronym{PICOC} framework
(Population, Intervention, Comparison, Outcome, Context) was utilized
\cite{kitchenham2007guidelines}. Details are presented in Table~\ref{tab:picoc}.

\begin{table}[ht]
	\fcmtcaption{Definition of the review scope using the PICOC framework.}\label{tab:picoc}
	\centering\footnotesize%
	\begin{tabular}{l >{\rightskip\fill}p{10cm}}
		\toprule
		Component             & Definition                                                                                                                                                               \\
		\midrule
		\textbf{P}opulation   & Deep Neural Networks (\acronym{DNN}) in Multi-Task Learning (\acronym{MTL}) settings, considering both hard and soft parameter sharing architectures.                    \\
		\textbf{I}ntervention & Adaptive Curriculum Learning (e.g., Self-Paced Learning, RL-based teachers) and Meta-Learning based loss reweighting (e.g., MAML, Learning to Reweight).                 \\
		\textbf{C}omparison   & 1. Static heuristics (Predefined curricula). \newline 2. Standard \acronym{MTL} baselines (Uniform weighting). \newline 3. Heuristic dynamic weighting (e.g., GradNorm). \\
		\textbf{O}utcome      & \textbf{Primary:} Mitigation of Negative Transfer, Generalization performance. \newline \textbf{Secondary:} Optimization stability, Convergence speed.                   \\
		\textbf{C}ontext      & Machine Learning research (excluding purely psychological curriculum studies without algorithmic implementation).                                                        \\
		\bottomrule
	\end{tabular}
\end{table}

\section{Research Questions (RQs)}

The following research questions have been defined to guide the data extraction
process:

\begin{description}
	\item[RQ1 (Curriculum)] How do adaptive approaches (e.g., Self-Paced Learning,
	      RL-based teachers) differ from static heuristics in defining training
	      schedules?
	\item[RQ2 (Multi-Task)] What strategies are employed to mitigate negative transfer
	      (e.g., gradient conflict resolution) in hard vs. soft parameter sharing
	      architectures?
	\item[RQ3 (Meta-Optimization)] How are meta-learning algorithms (e.g., MAML,
	      reweighting) applied to optimize loss weights or data selection dynamically?
\end{description}

\section{Search Strategy}

The search strategy relies on the intersection of three conceptual blocks:
\emph{Curriculum Learning}, \emph{Meta-Optimization}, and \emph{Multi-Task
	Learning}. Below are the specific queries designed for scientific databases.

\subsection{Scopus / Web of Science}

This is the primary source due to its broad indexing coverage.

\begin{verbatim}
TITLE-ABS-KEY (
  ("curriculum learning" OR "self-paced learning" OR "loss weighting"
   OR "adaptive weighting" OR "dynamic weighting")
  AND ("meta-learning" OR "meta-optimization" OR "learning to weight"
   OR "gradient-based" OR "differentiable" OR "automated")
  AND ("multi-task" OR "multi-objective" OR "auxiliary learning"
   OR "multi-loss" OR "negative transfer")
) AND ( DOCTYPE(ar) OR DOCTYPE(cp) ) AND NOT ( "reinforcement learning" )
\end{verbatim}

\subsection{IEEE Xplore}

A stricter syntax is required here (using the \texttt{"All Metadata"}
operator).

\begin{verbatim}
("All Metadata":"curriculum learning" OR "All Metadata":"self-paced learning"
 OR "All Metadata":"loss weighting" OR "All Metadata":"adaptive weighting"
 OR "All Metadata":"dynamic weighting")
AND ("All Metadata":"meta-learning" OR "All Metadata":"meta-optimization"
 OR "All Metadata":"learning to weight" OR "All Metadata":"gradient-based"
 OR "All Metadata":"differentiable" OR "All Metadata":"automated")
AND ("All Metadata":"multi-task" OR "All Metadata":"multi-objective"
 OR "All Metadata":"auxiliary learning" OR "All Metadata":"multi-loss")
\end{verbatim}

\subsection{Google Scholar}

Due to limitations in boolean logic handling, the search will be conducted
using shorter phrases:
\begin{verbatim}
"dynamic weighting" "multi task" "curriculum learning"
\end{verbatim}
Due to the large number of hits, only papers published after 2023 were
considered in the Google Scholar query.

\subsection{Snowballing Strategy}

To complement the database search and ensure comprehensive coverage, a
snowballing procedure will be applied \cite{wohlin2014guidelines}:
\begin{itemize}
	\item \textbf{Forward Snowballing:} Reviewing papers that cite seminal works identified in the supervisor's recommendations (e.g., \emph{GradNorm}, \emph{MAML}, \emph{Self-Paced Learning}) to identify recent advancements.
	\item \textbf{Backward Snowballing:} Examining the reference lists of included primary studies to identify foundational papers or missed relevant studies.
\end{itemize}

\section{Selection Criteria and Data Extraction}

The selection process will be conducted in two stages: title/abstract screening
followed by full-text review.

\begin{figure}[t]
	\centering
	\resizebox{0.95\textwidth}{!}{
		\begin{tikzpicture}[
    node distance=1.5cm,
    every node/.style={font=\small},
    block/.style={rectangle, draw, text width=6cm, align=center, minimum height=1.2cm},
    decision/.style={diamond, draw, text width=2cm, align=center, aspect=2},
    line/.style={draw, -latex'}
]

% Identification Phase
\node[block, text width=3.5cm] (scopus) {Scopus \\ (n = \textbf{17})};
\node[block, text width=3.5cm, right=0.4cm of scopus] (ieee) {IEEE Xplore \\ (n = \textbf{20})};
\node[block, text width=3.5cm, right=0.4cm of ieee] (scholar) {Google Scholar \\ (n = \textbf{69})};
\node[block, text width=3.5cm, right=0.4cm of scholar] (other) {Snowballing \\ (n = \textbf{18})};

% Screening Phase
\node[block, below=2cm of ieee, xshift=2cm] (duplicates) {Records after duplicates removed \\ (n = \textbf{123})};
\node[block, below=1cm of duplicates] (screened) {Records screened \\ (n = \textbf{56})};
\node[block, right=1cm of screened] (excluded) {Records excluded \\ (n = \textbf{67})};

% Eligibility Phase
\node[block, below=1cm of screened] (eligible) {Full-text articles assessed for eligibility \\ (n = \textbf{45})};
\node[block, right=1cm of eligible] (excluded_full) {Full-text articles excluded, with reasons \\ (n = \textbf{22})};

% Included Phase
\node[block, below=1cm of eligible] (included_qual) {Studies included in qualitative synthesis \\ (n = \textbf{23})};

% Connections
\path [line] (scopus) -- (duplicates);
\path [line] (ieee) -- (duplicates);
\path [line] (scholar) -- (duplicates);
\path [line] (other) -- (duplicates);
\path [line] (duplicates) -- (screened);
\path [line] (screened) -- (eligible);
\path [line] (screened) -- (excluded);
\path [line] (eligible) -- (included_qual);
\path [line] (eligible) -- (excluded_full);
% Removed connection to quantitative synthesis block

% Grouping labels (optional, can be added as nodes on the left)
\node[left=0.5cm of scopus, rotate=90, anchor=center] {\textbf{Identification}};
\node[left=0.5cm of screened, rotate=90, anchor=center] {\textbf{Screening}};
\node[left=0.5cm of eligible, rotate=90, anchor=center] {\textbf{Eligibility}};
\node[left=0.5cm of included_qual, rotate=90, anchor=center] {\textbf{Included}};

\end{tikzpicture}

	}
	\fcmfcaption{Schematic of the PRISMA selection process.}\label{rys:prisma}
\end{figure}

\paragraph{Inclusion Criteria.} The review will include papers that:
\begin{enumerate}
	\item Propose a method for \emph{automating} the training process (adaptive
	      curriculum or loss weighting).
	\item Address Multi-Task Learning challenges, specifically Negative Transfer or
	      parameter sharing.
	\item Utilize Meta-Learning, Self-Paced Learning, or Reinforcement Learning
	      approaches.
	\item Were published between 2017 and 2025.
\end{enumerate}

\paragraph{Data Extraction.} The following information will be extracted from selected papers:
\begin{itemize}
	\item \textbf{Metadata:} Title, year, publication venue.
	\item \textbf{Methodology:} Type of Curriculum (Static/Adaptive/SPL) and Optimization strategy.
	\item \textbf{Architecture:} Hard vs. Soft parameter sharing.
	\item \textbf{Results:} Mitigation of negative transfer and generalization improvements.
\end{itemize}
